\begin{helper}
Let $n,k\ge0$ be integers, let $p$ be real, let $A\colon\mathbb{Z}\to\mathbb{C}$
and let $\xi\in\mathbb{C}$ with $\xi\ne0$.  Write $A_k=A(k)$,
$\lambda_k=\lambda_k(n,p)$ for the exponent of \noderef{LambdaCoef},
$\epsilon_k=\epsilon_k(n,p)$ for the coefficient of \noderef{EpsCoef},
\[
  s=\lambda_k+n-k ,
\]
and let $T_k=T_k(n,p,A,\cdot)$ be the $k$-th summand of the series (19) as in
\noderef{Thm2Term}, so that
\[
  T_k(z)=\Bigl(A_kz^{k}+\epsilon_k\overline{A_k}\,\bar z^{\,k}\Bigr)\,|z|^{s}\,z^{-n}.
\]
Then, with $\partial/\partial\xi$ as in \noderef{DZ} and
$\partial/\partial\bar\xi$ as in \noderef{DZbar},
\[
  \frac{\partial T_k}{\partial\xi}(\xi)
  =\frac{|\xi|^{s}\,\xi^{-n}}{\xi}
    \Bigl(k\,A_k\xi^{k}
      +\Bigl(\frac{s}{2}-n\Bigr)
        \bigl(A_k\xi^{k}+\epsilon_k\overline{A_k}\,\bar\xi^{\,k}\bigr)\Bigr)
\]
and
\[
  \frac{\partial T_k}{\partial\bar\xi}(\xi)
  =\frac{|\xi|^{s}\,\xi^{-n}}{\bar\xi}
    \Bigl(k\,\epsilon_k\overline{A_k}\,\bar\xi^{\,k}
      +\frac{s}{2}
        \bigl(A_k\xi^{k}+\epsilon_k\overline{A_k}\,\bar\xi^{\,k}\bigr)\Bigr).
\]

These are the Cartesian form of the paper's polar transformation formulas
$H_{\bar\xi}=\tfrac12\bigl(H_r+\tfrac{i}{r}H_\theta\bigr)e^{i\theta}$ and
$H_\xi=\tfrac12\bigl(H_r-\tfrac{i}{r}H_\theta\bigr)e^{-i\theta}$ of
\texttt{paper.tex:234-237}, specialised to a single summand of (19); every
later use of equation (18) on one term reads them.

No constraint whatever is imposed on $n$, $p$ or $k$: the identity is a formal
differentiation of a product of integer and real powers, valid off the origin
for every exponent.  In particular neither $n\ge1$, nor $p>1$, nor $k\ge n$ is
needed.  That matters, because one consumer of this identity quantifies over
every integer $k\ge0$ with no lower bound at all.
\end{helper}
\begin{proof}
Throughout, $\xi\ne0$ is fixed, $s=\lambda_k+n-k$ is a fixed real number, and
$|z|^{s}$ means $\exp(s\log|z|)$ for $z\ne0$, read as a complex-valued
function of $z$.  Put
\[
  F(z)=A_kz^{k}+\epsilon_k\overline{A_k}\,\bar z^{\,k},\qquad
  G(z)=|z|^{s},\qquad
  H(z)=z^{-n},
\]
so that $T_k=F\cdot G\cdot H$ pointwise, by \noderef{Thm2Term}.

\emph{Step 0: two derivatives of the elementary functions.}  First, for the
identity function $\mathrm{id}(z)=z$ one has $D\,\mathrm{id}(z)=\mathrm{id}$,
so $D\,\mathrm{id}(z)[1]=1$ and $D\,\mathrm{id}(z)[i]=i$, whence by
\noderef{DZ} and \noderef{DZbar}
\begin{equation}\tag{1}
  \frac{\partial\,\mathrm{id}}{\partial\xi}=\frac12(1-i\cdot i)=1,
  \qquad
  \frac{\partial\,\mathrm{id}}{\partial\bar\xi}=\frac12(1+i\cdot i)=0 .
\end{equation}
Second, both operators are additive and commute with multiplication by a
complex constant on functions real-differentiable at the point in question,
because each is a fixed complex-linear expression in the values
$Df(\xi)[1]$ and $Df(\xi)[i]$ of the real Fr\'echet derivative, and
$D(f+g)(\xi)=Df(\xi)+Dg(\xi)$, $D(cf)(\xi)=c\,Df(\xi)$.

\emph{Step 1: smoothness of the three factors at $\xi$.}  The map
$z\mapsto\bar z$ is $\mathbb{R}$-linear, hence $C^\infty$ on $\mathbb{C}$; so
$z\mapsto z^{k}$ and $z\mapsto\bar z^{\,k}$ are $C^\infty$ on $\mathbb{C}$,
being integer powers of $C^\infty$ maps with $k\ge0$, and therefore $F$ is
$C^\infty$ on all of $\mathbb{C}$, in particular real-differentiable at $\xi$.
$G$ is $C^\infty$ at $\xi$ by part (a) of \noderef{WirtingerModulusPowerRule},
since $\xi\ne0$.  And $H$ is $C^\infty$ at $\xi$: on the open set
$\Omega=\mathbb{C}-\{0\}$ the function $z\mapsto z^{n}$ is $C^\infty$ and
nonvanishing, so its reciprocal $z\mapsto z^{-n}$ is $C^\infty$ there, and
$\xi\in\Omega$.

\emph{Step 2: the derivatives of $F$.}  Since $\overline{\,\cdot\,}$ is a ring
homomorphism, $\bar z^{\,k}=\overline{z^{k}}$.  Applying
\noderef{WirtingerPowerRule} with $\Omega=\mathbb{C}-\{0\}$, $g=\mathrm{id}$
and $a=k$, and then (1),
\begin{equation}\tag{2}
  \frac{\partial(z^{k})}{\partial\xi}(\xi)=k\,\xi^{k-1},
  \qquad
  \frac{\partial(z^{k})}{\partial\bar\xi}(\xi)=k\,\xi^{k-1}\cdot0=0 .
\end{equation}
Hence, by \noderef{WirtingerConjRule} applied to $f(z)=z^{k}$,
\begin{equation}\tag{3}
  \frac{\partial(\bar z^{\,k})}{\partial\xi}(\xi)
  =\overline{\frac{\partial(z^{k})}{\partial\bar\xi}(\xi)}=\bar0=0,
  \qquad
  \frac{\partial(\bar z^{\,k})}{\partial\bar\xi}(\xi)
  =\overline{\frac{\partial(z^{k})}{\partial\xi}(\xi)}
  =\overline{k\,\xi^{k-1}}=k\,\bar\xi^{\,k-1} .
\end{equation}
By the additivity and constant-multiplicativity of Step 0, together with (2)
and (3),
\begin{equation}\tag{4}
  \frac{\partial F}{\partial\xi}(\xi)=k\,A_k\,\xi^{k-1},
  \qquad
  \frac{\partial F}{\partial\bar\xi}(\xi)
  =k\,\epsilon_k\overline{A_k}\,\bar\xi^{\,k-1} .
\end{equation}

\emph{Step 3: the derivatives of $G$ and $H$.}  Part (b) of
\noderef{WirtingerModulusPowerRule} with $\mu=s$ gives
\begin{equation}\tag{5}
  \frac{\partial G}{\partial\xi}(\xi)=\frac{s}{2}\cdot\frac{|\xi|^{s}}{\xi},
  \qquad
  \frac{\partial G}{\partial\bar\xi}(\xi)
    =\frac{s}{2}\cdot\frac{|\xi|^{s}}{\bar\xi} .
\end{equation}
\noderef{WirtingerPowerRule} with $\Omega=\mathbb{C}-\{0\}$, $g=\mathrm{id}$
and $a=-n$, together with (1), gives
\begin{equation}\tag{6}
  \frac{\partial H}{\partial\xi}(\xi)=-n\,\xi^{-n-1},
  \qquad
  \frac{\partial H}{\partial\bar\xi}(\xi)=-n\,\xi^{-n-1}\cdot0=0 .
\end{equation}

\emph{Step 4: the product rule on the triple.}  All three factors are
real-differentiable at $\xi$ by Step 1, so two applications of
\noderef{WirtingerProductRule} --- first to $(FG)\cdot H$, then to $F\cdot G$
--- give, for each of the two operators $\partial\in\{\partial/\partial\xi,
\partial/\partial\bar\xi\}$,
\begin{equation}\tag{7}
  \partial(FGH)(\xi)
  =\bigl(\partial F(\xi)\,G(\xi)+F(\xi)\,\partial G(\xi)\bigr)H(\xi)
   +F(\xi)\,G(\xi)\,\partial H(\xi).
\end{equation}
(The product $FG$ is real-differentiable at $\xi$, being a product of two such
functions, so the outer application is legitimate.)

\emph{Step 5: collecting, for $\partial/\partial\xi$.}  Substituting (4), (5)
and (6) into (7),
\[
  \frac{\partial T_k}{\partial\xi}(\xi)
  =\Bigl(k\,A_k\xi^{k-1}\,|\xi|^{s}
     +F(\xi)\,\frac{s}{2}\frac{|\xi|^{s}}{\xi}\Bigr)\xi^{-n}
   +F(\xi)\,|\xi|^{s}\,\bigl(-n\,\xi^{-n-1}\bigr).
\]
Since $\xi\ne0$ we may write $\xi^{k-1}=\xi^{k}/\xi$ and
$\xi^{-n-1}=\xi^{-n}/\xi$, so every term carries the common factor
$|\xi|^{s}\xi^{-n}/\xi$, and factoring it out leaves
\[
  \frac{\partial T_k}{\partial\xi}(\xi)
  =\frac{|\xi|^{s}\xi^{-n}}{\xi}
    \Bigl(k\,A_k\xi^{k}+\frac{s}{2}F(\xi)-n\,F(\xi)\Bigr)
  =\frac{|\xi|^{s}\xi^{-n}}{\xi}
    \Bigl(k\,A_k\xi^{k}+\Bigl(\frac{s}{2}-n\Bigr)F(\xi)\Bigr),
\]
which, on substituting the definition of $F$, is the first asserted identity.

\emph{Step 6: collecting, for $\partial/\partial\bar\xi$.}  Substituting the
second halves of (4), (5) and (6) into (7), the last summand of (7) vanishes
because $\partial H/\partial\bar\xi(\xi)=0$, and
\[
  \frac{\partial T_k}{\partial\bar\xi}(\xi)
  =\Bigl(k\,\epsilon_k\overline{A_k}\,\bar\xi^{\,k-1}\,|\xi|^{s}
     +F(\xi)\,\frac{s}{2}\frac{|\xi|^{s}}{\bar\xi}\Bigr)\xi^{-n}.
\]
Since $\xi\ne0$ we also have $\bar\xi\ne0$, so
$\bar\xi^{\,k-1}=\bar\xi^{\,k}/\bar\xi$, and both terms carry the common
factor $|\xi|^{s}\xi^{-n}/\bar\xi$.  Factoring it out leaves
\[
  \frac{\partial T_k}{\partial\bar\xi}(\xi)
  =\frac{|\xi|^{s}\xi^{-n}}{\bar\xi}
    \Bigl(k\,\epsilon_k\overline{A_k}\,\bar\xi^{\,k}
      +\frac{s}{2}F(\xi)\Bigr),
\]
which, on substituting the definition of $F$, is the second asserted identity.

No step used any inequality between $n$, $k$ or $p$: the only hypothesis
consumed is $\xi\ne0$, needed for the smoothness of $G$ and $H$ at $\xi$, for
the three divisions by $\xi$ and $\bar\xi$, and for the negative powers
$\xi^{-n}$ to be the genuine reciprocals rather than the value $0$ assigned by
convention at the origin.
\end{proof}
